\scanpage{149}
$T^n(\alpha\beta)$ à l'aide de l'opération $*$ et des \deleted{\unsure{$\lambda$}} opérations
de restriction. Il faut voir comment s'expriment
les $T^n(\lambda^i(\alpha))$, on y reviendra.

Bien entendu, \deleted{\unsure{la} \illegible{1} \unsure{commutativité}} \inserted{les homomorphismes} $T_G^n$ sont fonctoriels :
si $u:G\longrightarrow G'$, on a \deleted{\unsure{cette}} commutativité \unsure{du|dans|des}
\diagram{\begin{tikzcd}[ampersand replacement=\&,column sep=large]
\text{(34)} \& A(G') \arrow[r,"T_{G'}"] \arrow[d,"A(u)"'] \& 1+\widehat{A(G')^*}{}^+ \arrow[d,"\widehat{A(u)}"]\\
\& A(G) \arrow[r,"T_G"] \& 1+\widehat{A(G)^*}{}^+
\end{tikzcd}}
% archive sheet number: 148
% printed form (vertical paper watermark, reverse face)
RITER \tikz[baseline=-.4ex,scale=.15]{\foreach \a in {0,45,...,315}{\draw (\a:.52) circle (.22);} \draw (0,0) circle (.21);} ARGENTINE \tikz[baseline=-.2ex,scale=.17]{\draw (0,0)--(1.6,0)--(1.8,.33)--(.5,.33)--(.32,.13)--(.92,.13);} \unsure{TYPEWRITE|TYPEWRITER}
